\section{Information sufficient for allocation}
A coordinator observes $Lw$ before choosing an assignment from a nonempty
finite set $D$. The unknown capability vector belongs to an unrestricted
finite-dimensional real vector space $W$; assignment $d$ gives disclosed
linear task reward $b_d(w)$. Observations are truthful, fixed linear
measurements. No strategic reporting, prior, positivity restriction or
computational bound is assumed. Fix $d_0$ and let $B$ have rows $b_d-b_{d_0}$.

\begin{theorem}[Allocation contrasts; AG-A1]
An $Lw$-based deterministic rule can choose an optimal assignment for every
$w\in W$ iff $\ker L\subseteq\ker B$. The minimum unrestricted scalar
linear measurement rank is $\operatorname{rank}B$.
\end{theorem}
\begin{proof}
For $t\in\ker L$, the assignment chosen at zero data must maximize reward
at both $t$ and $-t$. Every reward functional consequently agrees on $t$,
so $Bt=0$. Conversely the kernel inclusion makes all reward contrasts
well-defined functions of $Lw$, which suffice for optimization. Rank-nullity
gives the lower bound; measuring $Bw$ attains it.
\end{proof}
This is a standard linear-estimability argument applied to decisions, not a
new general theory of decision-focused prediction \citep{elmachtoub2020}.
For a restricted capability domain $C$, the correct condition is a common
maximizer on each observation fiber:
\[
 \bigcap_{w\in C:\,Lw=y}\arg\max_{d\in D}b_d(w)\ne\varnothing.
\]
For example, $b_0=0,b_1(w)=w$ on $w\ge0$ needs no measurements, unlike
the full-line problem. Offsets and bounded domains cannot be ignored.

For $n$ agents filling $n$ fixed tasks bijectively, write $C_{ij}$ for arbitrary
real completion scores. Assignment $\pi$ has score $\sum_i C_{i,\pi(i)}$.
The reward-contrast span consists of matrices with zero row and column sums,
with basis
$E_{ij}-E_{in}-E_{nj}+E_{nn}$, $i,j<n$. Each basis rectangle is a difference
of two assignments exchanging two task choices. Thus exactly $(n-1)^2$
unrestricted scalar linear measurements suffice and are necessary.

By comparison, testing physical capability-row interchangeability of two agents
with tasks fixed has discrepancy $T(C)=(C_{1j}-C_{2j})_j$, rank $n$.
This is a capability-array test, not a claim about a complete normal-form game's
automorphism group. The canonical Game Fairness criterion identifies $Tw=0$
through linear observations exactly when their kernel is contained in $\ker T$.
For $n=2$ allocation rank $1$ is smaller than test rank $2$; for $n=3$
allocation rank $4$ is larger than test rank $3$. Their row spaces at $n=3$
are incomparable: a constant added to one row changes interchangeability but
not any assignment comparison; changing only $C_{31}$ from positive to
negative leaves the first two rows equal but changes which tasks can be
optimal. The joint measurement rank is $5$, since the row-space intersection
has dimension $2$.

\paragraph{Noise, messages and incentives.}
If reward contrasts factor as $b_d-b_{d_0}=a_dL$ and observations have error
$\|e\|\le\epsilon$, maximizing $a_d\cdot(Lw+e)$ incurs regret at most
\[
 \epsilon\max_{d,e}\|a_d-a_e\|_*.
\]
Observed optimality and the dual-norm inequality prove this; for two decisions
it is the sharp deterministic minimax value on a full linear class with
surjective $L$. It is not a randomized-regret or finite-bit bound. An informed
sender can instead communicate a selected assignment using $n!$ messages,
which is sharp since each permutation is uniquely optimal somewhere. Nor does
truthful measurement imply truthful reports: a low-capability agent paid only
for receiving a task may profitably exaggerate its capability. Full proofs
and exact witnesses remain in the allocation research note.
